%=============================================================================!
% 8.0  CHAPTER OPENING                                                        !
%=============================================================================!
The preceding chapters treated the potential energy $U(\vb{r}^N)$ as a
known function and used its derivatives to determine the motion. For an
atomistic simulation we must also choose how to obtain that function.
Within the Born--Oppenheimer approximation, the electrons remain in their
ground state as the nuclei move, and their energy together with the
nuclear repulsion defines the potential energy surface. Evaluating an
electronic-structure calculation at every step is costly, which motivates
simpler interatomic models.

We begin by examining when classical motion of the nuclei is an
appropriate approximation. Pair potentials then give a direct route from
an energy model to forces, their numerical checks and the treatment of a
cutoff. Their limitations motivate environment-dependent interactions for
metals and carbon, and bonded force fields for molecules. Finally, the
slow decay of the Coulomb interaction shows why charges need a separate
summation method. Chapter~9 uses these forces to advance the motion,
while Chapter~18 introduces the learning methods used to represent a
potential energy surface. Later chapters on these models are planned
for the continuation of Part~V.
